% ============================================================
\section{社会福利函数：从个人效用加总成社会评价（补全）}
% ============================================================

政策评估需要把 $N$ 个人的效用（或收入）映射为一个数——
这个映射本身就是一类“效用函数”，其形式选择\textbf{直接编码分配价值观}。

\subsection{功利主义（效用加总型）}

\notehead{形式}
$W=\sum_i u_i$。

\notehead{性质}
每个人的效用权重相同；对分配\textbf{完全中性}：
一元从穷人转给富人，$W$ 不变（若边际效用不被考虑）；
若 $u$ 凹，则同样总量下平均分配最优——\textbf{分配关怀全部来自个体效用曲率，
而非 $W$ 的形式}。
优点：与补偿原则（总剩余最大化）一致，是成本收益分析的默认件；
缺点：无法接受“一元对穷人的价值更高”这类独立的分配判断。

\begin{figure}[H]
\centering
\begin{tikzpicture}[font=\scriptsize,scale=0.88,transform shape]
  \def\frontier{(0,0.6) .. controls (2.4,3.3) and (1.5,3.5) .. (3.3,3.6)}
  % (a) 功利主义
  \draw[axisline] (0,0) -- (3.45,0) node[below right] {$u_1$};
  \draw[axisline] (0,0) -- (0,3.9) node[above left] {$u_2$};
  \draw[thick] \frontier;
  \draw[dashed,thin] (0,3.6) -- (3.3,3.6);
  \draw[->,thick] (0,3.6) -- (3.3,0) node[midway,above left,font=\scriptsize] {$u_1+u_2=\text{常数}$};
  \fill (1.95,1.65) circle (1.5pt);
  \node[font=\scriptsize, align=center] at (1.95,1.20) {(a) 功利主义：\\最大化总和};
  % (b) 最大最小
  \begin{scope}[xshift=5.3cm]
    \draw[axisline] (0,0) -- (3.45,0) node[below right] {$u_1$};
    \draw[axisline] (0,0) -- (0,3.9) node[above left] {$u_2$};
    \draw[thick] (0,0.6) .. controls (2.4,3.3) and (1.5,3.5) .. (3.3,3.6);
    \draw[dashed,thin] (0.0,2.78) -- (2.42,2.78);
    \draw[->,thick] (0,2.78) -- (2.58,0) node[midway,above left,font=\scriptsize] {$u_2=\text{常数}$};
    \fill (2.42,2.78) circle (1.5pt);
    \node[align=center,font=\scriptsize] at (2.05,1.35) {(b) 最大最小：\\抬高最差者的水平};
  \end{scope}
  % (c) CES
  \begin{scope}[xshift=10.6cm]
    \draw[axisline] (0,0) -- (3.45,0) node[below right] {$u_1$};
    \draw[axisline] (0,0) -- (0,3.9) node[above left] {$u_2$};
    \draw[thick] (0,0.6) .. controls (2.4,3.3) and (1.5,3.5) .. (3.3,3.6);
    \draw[->,thick] (0,2.20) to[bend left=22] (2.95,0.66) node[midway,above,font=\scriptsize] {等 $\varepsilon$ 无差异曲线};
    \fill (2.16,1.20) circle (1.5pt);
    \node[align=center,font=\scriptsize] at (2.05,2.42) {(c) CES：$\varepsilon$ 居中，\\切于两极限之间};
  \end{scope}
\end{tikzpicture}
\caption{三种社会福利函数在同一可能性前沿上的最优分配}
\end{figure}
\figexplan{三幅图共用同一条效用可能性前沿（实线），区别只在等社会福利线。功利主义用斜率为 $-1$ 的直线（$u_1+u_2$ 为常数），最优点在两效用较均衡处；最大最小用水平线（$u_2$ 为常数），最优点是把最差者的效用推到最高；CES 型用一族曲率居中的无差异曲线，切点介于前两者之间。图形要点：福利函数的选择等价于在可能性前沿上“选点”，参数 $\varepsilon$ 正是这一选择的调节旋钮。}

\subsection{最大最小（最差者优先型）}

\notehead{形式}
$W=\min_i u_i$。

\notehead{性质}
社会福利等于\textbf{最差者}的效用；任何不改善最差者的再分配都无价值。
无差异曲线为 L 形（与个人效用的完全互补同构，§4.1.3）：
只有最差者改善，$W$ 才改善。
极限性质：它对不平等“无限厌恶”——宁可总量大幅缩水也要托底。

\subsection{CES 型社会福利函数（含不平等厌恶参数）}

\notehead{形式}
\[
W(x)=\Bigl[\frac1N\sum_i x_i^{\,1-\varepsilon}\Bigr]^{\frac{1}{1-\varepsilon}},
\qquad \varepsilon\ge0,\ \varepsilon\ne1,
\]
$\varepsilon$ 为\textbf{不平等厌恶}参数；$\varepsilon=1$ 时取几何平均
$\prod_i x_i^{1/N}$。

\notehead{两个极限（推导）}
$\varepsilon\to0$：$W\to\frac1N\sum_i x_i$——功利主义（算术平均）；
$\varepsilon\to\infty$：幂均值的极限为最小值，$W\to\min_i x_i$——最大最小。
故 CES 型是\textbf{功利主义与最大最小之间的连续桥梁}，
$\varepsilon$ 就是这座桥上的刻度。

\notehead{均等分布等价与阿特金森指数}
定义\textbf{均等分布等价}（EDE）$\xi=W(x)$：
若每人得 $\xi$，社会福利与现实分配相同。不平等指数
\[
A=1-\frac{\xi}{\bar x},\qquad \bar x=\frac1N\sum_ix_i.
\]
\textbf{数值例}：两人收入 $(10,2)$，$\bar x=6$。
\begin{center}
\small
\begin{tabular}{lccc}
\toprule
$\varepsilon$ & $\xi$（计算） & $A=1-\xi/6$ & 含义 \\
\midrule
$0$（功利主义） & $6.0000$ & $0.0000$ & 只看总量 \\
$1$（几何平均） & $\sqrt{20}=4.4721$ & $0.2546$ & 温和厌恶 \\
$2$ & $(0.5(10^{-1}+2^{-1}))^{-1}=3.3333$ & $0.4444$ & 强厌恶 \\
$\to\infty$（最大最小） & $\to2.0000$ & $\to0.6667$ & 只看最差者 \\
\bottomrule
\end{tabular}
\end{center}
（复算 $\varepsilon=2$：$W=[0.5(10^{-1}+2^{-1})]^{-1}=[0.5\times0.6]^{-1}=1/0.3=3.3333$。）

\begin{center}
\begin{tikzpicture}[scale=0.95]
  \draw[axisline] (0,0) -- (5.9,0) node[below left, axislabel] {$u_1$};
  \draw[axisline] (0,0) -- (0,4.4) node[below left, axislabel] {$u_2$};
  % epsilon=0: straight line
  \draw[thick, MidBlue] (0,3.4) -- (4.4,1.0);
  \node[right] at (4.4,1.05) {$\varepsilon=0$};
  % epsilon=1: hyperbola
  \draw[thick, MidBlue, domain=0.8:5.0, smooth, variable=\x] plot ({\x},{2.2/\x});
  \node[above right] at (3.55,1.05) {$\varepsilon=1$};
  % epsilon large: near L
  \draw[thick, MidBlue] (0.95,3.8) .. controls (1.0,1.45) and (1.4,1.0) .. (4.6,0.75);
  \node[right] at (4.6,0.62) {$\varepsilon$ 大};
  \node[align=left] at (2.9,3.8) {社会无差异曲线：\\ 曲率 $=$ 不平等厌恶程度\\ L 形极限 $=$ 最大最小};
\end{tikzpicture}
\end{center}
\figexplan{横纵轴为两人的效用。$\varepsilon=0$ 时为直线（功利主义，只关心总和）；$\varepsilon=1$ 时为双曲线（乘法型）；$\varepsilon$ 越大越接近 L 形，极限即最大最小（罗尔斯）准则。曲率因此可直接读作不平等厌恶程度：政策评估只需随 $\varepsilon$ 重算一遍。}

\notehead{为什么选它}
(i) 一参数覆盖全部主流分配立场，政策报告的\textbf{分配敏感性}
就是随 $\varepsilon$ 重算一遍；
(ii) EDE 与不平等指数有闭式，沟通成本低；
(iii) 与个体 CES 同构，宏观模型中“社会规划者偏好”可直接借用消费品加总技术。

\subsection{Kolm--Pollak 型：绝对不等价}

\notehead{形式}
\[
W(x)=-\frac{1}{\kappa}\ln\Bigl(\frac1N\sum_i e^{-\kappa x_i}\Bigr),
\qquad \kappa>0,
\]
\textbf{性质}：给所有人同时增加等量收入，$W$ 增加相同的量
（平移不变性）——不平等按其\textbf{绝对差距}度量；
CES 型则按\textbf{比例差距}（$x$ 与 $2x$ 分布被认为同样不平等）。
选择标准：认为“贫富差距一万”与“贫富差距十万，但同比例”一样不平等
→ Kolm--Pollak；认为不同 → CES 型。
$W$ 同时是\textbf{确定性等价}结构：与 CARA 同族，
税率/转移的福利分析中常与 CRRA 个体效用搭配使用。

\subsection{几何平均与议价型福利函数}

\notehead{形式}
$W=\prod_i u_i^{1/N}$（CES 型 $\varepsilon=1$ 特例），
等价于 $\ln W=\frac1N\sum_i\ln u_i$。
两大用途：
(i) 两人、可变权重时即为议价解的目标函数（§4.2.6），
\textbf{社会选择与合作解在此汇合}；
(ii) 代际福利：$W=\prod_t c_t^{\beta_t}$ 对不完全替代的不同时期消费
赋予零下限（任何一期消费为零，$W=0$）——
气候经济学中“不允许未来被完全牺牲”的形式化方式之一。

\subsection{社会福利函数选择对照表}

\begin{center}
\small
\begin{tabular}{llll}
\toprule
形式 & 不平等态度 & 关键性质 & 典型用途 \\
\midrule
功利主义 $\sum u_i$ & 中性 & 与补偿原则一致 & 成本收益分析 \\
最大最小 $\min u_i$ & 无限厌恶 & 只认最差者 & 托底政策、正义论 \\
CES 型（$\varepsilon$） & 一参数可调 & 比例差距、EDE 闭式 & 分配敏感性报告 \\
Kolm--Pollak（$\kappa$） & 一参数可调 & 绝对差距、平移不变 & 税与转移分析 \\
几何平均 & 温和厌恶 & 与议价解同构 & 代际福利 \\
\bottomrule
\end{tabular}
\end{center}
